\chapter{Conjunctions and Word Order}

German connectors do not all affect word order in the same way. The most
useful distinction is between coordinating conjunctions, subordinating
conjunctions, and connecting adverbs.

\begin{tcolorbox}[
  colback=black!3,
  colframe=black!55,
  title={Three patterns to remember},
  fonttitle=\bfseries
]
\begin{tabularx}{\textwidth}{@{}>{\bfseries}l X@{}}
Hauptsatz & subject $+$ \textbf{finite verb} $+$ remaining information \\
Nebensatz & conjunction $+$ subject $+$ remaining information $+$ \textbf{finite verb} \\
Connector adverb & adverb $+$ \textbf{finite verb} $+$ subject $+$ remaining information
\end{tabularx}
\end{tcolorbox}

\section{Coordinating conjunctions: Hauptsatz + Hauptsatz}

Coordinating conjunctions connect elements of equal rank. When they connect
two main clauses, they occupy ``position zero'': the normal verb-second word
order of the following Hauptsatz does not change.

\begin{center}
\small
\rowcolors{2}{referenceBlueBg}{white}
\begin{tabularx}{\textwidth}{@{}>{\bfseries}l l l X@{}}
\toprule
German & English & Function & Example \\
\midrule
und & and & addition & Ich koche, und er \textbf{deckt} den Tisch. \\
oder & or & alternative & Wir fahren heute, oder wir \textbf{fahren} morgen. \\
aber & but & contrast & Sie ist müde, aber sie \textbf{arbeitet} weiter. \\
denn & because/for & reason & Ich bleibe zu Hause, denn ich \textbf{bin} krank. \\
sondern & but rather & correction after a negative & Er fährt nicht, sondern er \textbf{nimmt} den Zug. \\
sowie & as well as & addition & Sie spricht Deutsch sowie Englisch. \\
\bottomrule
\end{tabularx}
\end{center}

\noindent
\textbf{Pattern:}
\quad Hauptsatz, $+$ conjunction $+$ subject $+$ \textbf{finite verb} $+$ \ldots

\section{Subordinating conjunctions: Hauptsatz + Nebensatz}

A subordinating conjunction introduces a dependent clause. The finite verb
moves to the end of that Nebensatz, and a comma separates it from the main
clause.

\subsection{Content, reason, condition, and contrast}

\begin{center}
\small
\rowcolors{2}{referenceGreenBg}{white}
\begin{tabularx}{\textwidth}{@{}>{\bfseries}l l l X@{}}
\toprule
German & English & Function & Example \\
\midrule
dass & that & statement/content & Ich weiß, dass er heute \textbf{kommt}. \\
ob & whether/if & indirect yes/no question & Sie fragt, ob der Laden offen \textbf{ist}. \\
weil & because & reason & Wir gehen nicht, weil es stark \textbf{regnet}. \\
da & since/because & reason, often already known & Da es spät \textbf{ist}, gehen wir nach Hause. \\
wenn & if/when & condition; repeated/present/future time & Ich komme, wenn ich Zeit \textbf{habe}. \\
falls & if/in case & possible condition & Ruf mich an, falls du Hilfe \textbf{brauchst}. \\
obwohl & although & concession/contrast & Er geht spazieren, obwohl es \textbf{regnet}. \\
während & while/whereas & simultaneous time or contrast & Sie liest, während er Musik \textbf{hört}. \\
\bottomrule
\end{tabularx}
\end{center}

\subsection{Time, purpose, result, and manner}

\begin{center}
\small
\rowcolors{2}{referenceGreenBg}{white}
\begin{tabularx}{\textwidth}{@{}>{\bfseries}l l l X@{}}
\toprule
German & English & Function & Example \\
\midrule
als & when & one event/period in the past & Als ich klein \textbf{war}, wohnte ich in Bonn. \\
bevor/ehe & before & earlier event & Wasch dir die Hände, bevor du \textbf{isst}. \\
nachdem & after & earlier completed event & Nachdem wir gegessen \textbf{hatten}, gingen wir spazieren. \\
seit/seitdem & since & starting point continuing until now & Seitdem sie hier \textbf{wohnt}, spricht sie mehr Deutsch. \\
bis & until & end point & Warte, bis der Bus \textbf{kommt}. \\
sobald & as soon as & immediate sequence & Ich rufe an, sobald ich zu Hause \textbf{bin}. \\
solange & as long as & duration/condition & Bleib hier, solange du \textbf{möchtest}. \\
damit & so that & purpose & Ich spreche langsam, damit alle mich \textbf{verstehen}. \\
sodass/so dass & so that/as a result & consequence & Es schneite, sodass die Straßen glatt \textbf{wurden}. \\
indem & by/by means of & manner or method & Man lernt, indem man regelmäßig \textbf{übt}. \\
ohne dass & without & absent accompanying action & Er ging, ohne dass er sich \textbf{verabschiedete}. \\
\bottomrule
\end{tabularx}
\end{center}

\subsection{When the Nebensatz comes first}

The complete Nebensatz occupies position one of the following Hauptsatz.
Consequently, the finite verb of the Hauptsatz comes immediately after the
comma.

\begin{center}
\rowcolors{2}{referenceGreenBg}{white}
\begin{tabularx}{\textwidth}{@{}l X@{}}
\toprule
Order & Example \\
\midrule
HS + NS & Ich bleibe heute zu Hause, weil ich krank \textbf{bin}. \\
NS + HS & Weil ich krank \textbf{bin}, \textbf{bleibe} ich heute zu Hause. \\
\bottomrule
\end{tabularx}
\end{center}

\section{Connecting adverbs: Hauptsatz with inversion}

Words such as \textit{deshalb} and \textit{trotzdem} are adverbs rather than
conjunctions. The following clause remains a Hauptsatz. When the adverb stands
in position one, the finite verb must stand directly after it in position two,
before the subject.

\begin{center}
\small
\rowcolors{2}{referenceYellowBg}{white}
\begin{tabularx}{\textwidth}{@{}>{\bfseries}l l l X@{}}
\toprule
German & English & Function & Example \\
\midrule
deshalb/deswegen & therefore/for that reason & result & Ich bin krank; deshalb \textbf{bleibe} ich zu Hause. \\
trotzdem & nevertheless/even so & contrast & Es regnet; trotzdem \textbf{gehen} wir spazieren. \\
außerdem & moreover/in addition & addition & Sie spricht Deutsch; außerdem \textbf{lernt} sie Spanisch. \\
dann & then & sequence/result & Wir essen; dann \textbf{gehen} wir ins Kino. \\
danach & afterwards & sequence & Der Kurs endet; danach \textbf{fahre} ich nach Hause. \\
sonst & otherwise & alternative/consequence & Beeil dich; sonst \textbf{verpassen} wir den Bus. \\
jedoch & however & contrast & Er war müde; jedoch \textbf{arbeitete} er weiter. \\
\bottomrule
\end{tabularx}
\end{center}

\noindent
\textbf{Compare:}
\quad \textit{weil ich krank \textbf{bin}} (Nebensatz)
\quad but \quad
\textit{deshalb \textbf{bleibe} ich zu Hause} (Hauptsatz).

\section{Paired conjunctions}

\begin{center}
\small
\rowcolors{2}{referenceRedBg}{white}
\begin{tabularx}{\textwidth}{@{}>{\bfseries}l l X@{}}
\toprule
German & English & Example \\
\midrule
sowohl \ldots{} als auch & both \ldots{} and & Sie spricht sowohl Deutsch als auch Englisch. \\
entweder \ldots{} oder & either \ldots{} or & Entweder kochen wir, oder wir bestellen eine Pizza. \\
weder \ldots{} noch & neither \ldots{} nor & Er trinkt weder Kaffee noch Tee. \\
nicht nur \ldots{} sondern auch & not only \ldots{} but also & Sie ist nicht nur klug, sondern auch geduldig. \\
zwar \ldots{} aber & admittedly \ldots{} but & Der Weg ist zwar lang, aber er ist schön. \\
je \ldots{} desto/umso & the \ldots{} the & Je mehr du übst, desto leichter \textbf{wird} es. \\
\bottomrule
\end{tabularx}
\end{center}

With \textit{je \ldots{} desto/umso}, the \textit{je}-clause has subordinate
word order, while the \textit{desto/umso}-clause has main-clause word order:
\textit{Je mehr du \textbf{übst}, desto leichter \textbf{wird} es.}

\section{Quick decision guide}

\begin{enumerate}[leftmargin=*]
  \item If the connector is \textit{und, oder, aber, denn}, or \textit{sondern},
        keep normal Hauptsatz word order after it.
  \item If it is a subordinating conjunction such as \textit{weil, dass,
        obwohl}, or \textit{wenn}, put the finite verb at the end of its clause.
  \item If it is a connecting adverb such as \textit{deshalb, trotzdem}, or
        \textit{außerdem}, put the finite verb immediately after the connector.
  \item Separate a Nebensatz from its Hauptsatz with a comma.
\end{enumerate}
